\setClassNameDisplay{SongInPlaylist}{Song\_in\_Playlist}
\ifbeamer
    \definecolor{relColor}{RGB}{0, 210, 255}
\else
    \definecolor{relColor}{RGB}{0, 140, 180}
\fi

\Aufgabe{Song-Datenbank Diagramme: \texorpdfstring{\url{www.dbiu.de/songs}}{www.dbiu.de/songs}}{

Zeichne die \emphColB{Klassenkarten} der Tabellen \emphColB{Song} und \emphColB{Playlist}. Zeichne anschließend mit \emphColC{zwei verschiedenen Farben} die beiden \emphColC{Darstellungsmöglichkeiten der Beziehung} zwischen den beiden Tabellen ein.

\LoesungKaro{
\begin{center}
\begin{tikzpicture}[scale=0.78, transform shape]
    \begin{class}[text width=4.2cm]{Song}{0,0}
        \attribute{int id}
        \attribute{String titel}
        \attribute{String genre}
        \attribute{String interpret}
    \end{class}

    \begin{class}[text width=3.8cm]{Playlist}{11,0}
        \attribute{int id}
        \attribute{String titel}
    \end{class}

    % Farbe 2: Beziehungstabelle in relColor
    {
    \renewcommand{\umldrawcolor}{relColor}
    \renewcommand{\umltextcolor}{relColor}
    \begin{class}[text width=4.5cm]{SongInPlaylist}{5.5,-2.6}
        \attribute{int song\_id}
        \attribute{int playlist\_id}
    \end{class}
    }

    % Farbe 1: direkte m:n-Beziehung
    \draw[very thick, colA] 
        ([shift={(0,0.5)}]Song.east) to 
            node[pos=0.08, above=2pt]{m}
            node[midway, above=2pt]{\emph{enthalten\_in}}
            node[pos=0.92, above=2pt]{n}
        ([shift={(0,0.5)}]Playlist.west);

    % Farbe 2: mit Beziehungstabelle
    \draw[very thick, -{Straight Barb[scale=1.5]}, relColor] 
        ([shift={(0,0.3)}]SongInPlaylist.west) to[out=180, in=270] 
            node[pos=0.1, above=2pt]{n}
            node[midway, below left]{song\_id}
            node[pos=0.9, left=2pt]{1}
        (Song.south);

    \draw[very thick, -{Straight Barb[scale=1.5]}, relColor] 
        ([shift={(0,0.3)}]SongInPlaylist.east) to[out=0, in=270] 
            node[pos=0.1, above=2pt]{n}
            node[midway, below right]{playlist\_id}
            node[pos=0.9, right=2pt]{1}
        (Playlist.south);

    % Legende
    \node[anchor=north, align=center] at (5.5,-4.8) {
        \textbf{\textcolor{colA}{Farbe 1: Direkte m:n-Beziehung}}\qquad\qquad
        \textbf{\textcolor{relColor}{Farbe 2: Mit Beziehungstabelle (zwei 1:n-Beziehungen)}}
    };
\end{tikzpicture}
\end{center}
}{15}
}